\section*{Hoofdstuk \ref{ch:b2:potential-wells-tunneling} ---
Potentiaalputten, barrières en tunneling}

\begin{solution}{exo:b2:potential-wells-tunneling:1}
\qty{1}{nm}: \qty{0.38}{}, \qty{1.50}{}, \qty{3.38}{eV}; $2 \to 1$:
\qty{1.13}{eV}, \qty{1.1}{\micro m}.
\qty{0.1}{nm}: \qty{38}{eV}. Nucleon: \qty{8}{MeV}. Knikker:
$E_1 = \qty{5.5e-63}{J}$;
$n = 2mvL/h = 3 \times 10^{27}$.
\end{solution}

\begin{solution}{exo:b2:potential-wells-tunneling:2}
Opsluiting $\qty{3.76}{eV.nm^2}/L^2$: $0.10$, $0.24$, $0.42$, \qty{0.60}{eV}; totaal
$1.84$, $1.98$, $2.16$, \qty{2.34}{eV}: \qty{673}{}, \qty{628}{},
\qty{574}{}, \qty{530}{nm} ---
rood, oranje, geelgroen, groen.
\end{solution}

\begin{solution}{exo:b2:potential-wells-tunneling:3}
$\kappa = \qty{1.09e10}{m^{-1}}$; $\eu^{-2\kappa a}$: $1.5 \times 10^{-3}$,
$1.8 \times 10^{-5}$,
$3 \times 10^{-10}$, $10^{-19}$; $\times 8.8$ per \qty{0.1}{nm}; proton: $\kappa$ is
$43$ keer groter, $\eu^{-280}$ bij \qty{0.3}{nm}.
\end{solution}

\begin{solution}{exo:b2:potential-wells-tunneling:4}
$k_1 = \qty{7.3e9}{m^{-1}}$, $k_2 = \qty{5.1e9}{m^{-1}}$; $R = 0.03$,
$T = 0.97$. Trap
omlaag: $k_2 = \qty{8.9e9}{m^{-1}}$, $R = 0.01$. Een klassiek deeltje keert nooit
terug wanneer het de energie heeft; een lichtgolf wordt bij elke verandering van
index gedeeltelijk weerkaatst.
\end{solution}

\begin{solution}{exo:b2:potential-wells-tunneling:5}
(a), (b) \cref{prop:b2:potential-wells-tunneling:infinite}. (c) $L/2$;
$0.18L$; $0$; $\pi\hbar/L$; product $0.57\hbar > \hbar/2$. (d) Snelle
trillingen rond het gemiddelde $1/L$ --- de klassieke gelijkmatige dichtheid.
\end{solution}

\begin{solution}{exo:b2:potential-wells-tunneling:6}
(a) \cref{prop:b2:potential-wells-tunneling:finite}. (b) $1$, $3$, $7$.
(c) $R = 2.6$, twee toestanden; $\kappa = \qty{4.4e9}{m^{-1}}$, diepte
\qty{0.23}{nm}. (d)
Oneindig veel, van onderaf naar die van de oneindige put toe.
\end{solution}

\begin{solution}{exo:b2:potential-wells-tunneling:7}
(a) $1 + r = t$, $\iu k(1 - r) = -\kappa t$. (b) Teller en noemer
zijn toegevoegd complex; fase $-2\arctan(\kappa/k)$. (c) $\psi^*\psi'$ is
reëel voor een
reële exponentiële. (d) \qty{0.2}{nm}; zoals de \omterm{def:b2:dispersion-wave-packets:complex}{evanescente golf} van de totale
weerkaatsing: aanwezig, zonder flux te dragen, en in staat een tweede midden
te voeden
--- dat is tunnelen.
\end{solution}

\begin{solution}{exo:b2:potential-wells-tunneling:8}
(a) \cref{thm:b2:potential-wells-tunneling:tunnel}. (b) $\kappa a = 3$:
exact $1/(1 + \sinh^23) = 9.9 \times 10^{-3}$, benadering
$4\eu^{-6} = 9.9 \times
10^{-3}$; $\kappa a = 1$: $0.42$ tegen $0.54$. (c)
$\kappa \to \iu k_2$; $T = 1$
voor $k_2a = n\pi$. (d) De weerkaatsingen van de twee vlakken heffen elkaar
op wanneer de
barrière een geheel aantal halve golflengten is --- resonante
doorlating.
\end{solution}

\begin{solution}{exo:b2:potential-wells-tunneling:9}
(a) $b = \qty{52}{fm}$; $G = 82$; $T = \eu^{-82} \approx 2 \times 10^{-36}$.
(b) $v =
\qty{1.55e7}{m/s}$, $10^{21}$ botsingen per seconde: tempo
$2 \times 10^{-15}\,\unit{s^{-1}}$,
levensduur $\sim 10^7$ jaar. (c) $G = 100$ en $69$: $\eu^{18} \approx 10^8$ langer,
$\eu^{-13} \approx 10^{-6}$ korter. (d) Het $\alpha$-deeltje is stevig
gebonden en vertrekt
met positieve energie; een proton niet; een $^{12}$C heeft drie keer de
lading en een veel grotere exponent (clusterverval bestaat, bij $10^{-10}$).
\end{solution}

\begin{solution}{exo:b2:potential-wells-tunneling:10}
(a) \qty{24}{GHz}; \qty{41}{ps}. (b)
$\psi = (\varphi_+\eu^{-\iu E_+t/\hbar} +
\varphi_-\eu^{-\iu E_-t/\hbar})/\sqrt2$;
$P = \sin^2(\pi\nu t)$. (c) $\eu^{(\sqrt2 - 1)
\kappa a} = 15$:
$\kappa a = 6.5$. (d) Zwaardere groepen en hogere barrières maken
de tunnelperiode langer dan jaren.
\end{solution}

\begin{solution}{exo:b2:potential-wells-tunneling:11}
(a) \qty{1.09e10}{m^{-1}}; $\times 8.8$; $1\%$ van $I$ is \qty{0.5}{pm}. (b)
$6 \times 10^9$.
(c) Elke buur op $d' = \qty{0.56}{nm}$ draagt
$\eu^{-2\kappa \times 0.06\,\text{nm}}
= 0.28$ van de stroom van de top; de
top draagt ongeveer de helft. (d)
$10^{-5}\,\unit{K}$ --- onmogelijk; symmetrische ontwerpen waarin punt en monster
samen uitzetten, snelle aftastingen, materialen met lage uitzetting, en de
\omterm{def:b2:feedback-oscillators:loop}{terugkoppeling}.
\end{solution}

\begin{solution}{exo:b2:potential-wells-tunneling:12}
(a) \qty{0.056}{} en \qty{0.008}{eV}. (b) \qty{1.48}{eV}, \qty{835}{nm}. (c)
Opsluiting
\qty{0.17}{eV}: $L = \qty{6.1}{nm}$. (d) $R = 3.6$: drie.
\end{solution}

\begin{solution}{pb:b2:potential-wells-tunneling:1}
\textbf{1.} $\kappa = \sqrt{2m\Phi}/\hbar = \qty{1.09e10}{m^{-1}}$.

\textbf{2.} $\eu^{2.2} = 8.8$; $\eu^{0.022}$: $2.2\%$ per picometer.

\textbf{3.} \qty{0.5}{pm}.

\textbf{4.} $6 \times 10^9$ per seconde; \qty{0.16}{ns} uit elkaar, tegen
$\hbar/\Phi
\sim 10^{-16}\,\unit{s}$: één tegelijk.

\textbf{5.} Ongeveer de helft; de rest dooft zo snel met de afstand uit dat het
beeld dat van één atoom is.

\textbf{6.} $\times 8.8$; vaste stroom houdt de punt veilig en maakt van
de exponentiële een lineaire hoogte.

\textbf{7.} $\Delta T = 10^{-12}/(10^{-5} \times 0.02) = \qty{5e-6}{K}$;
symmetrische
ophangingen, snelheid, materialen met lage uitzetting, en een \omterm{def:b2:feedback-oscillators:loop}{terugkoppeling} die de
drift volgt.

\textbf{8.} Een trapezium, verlaagd aan de opvangende kant: de werkzame
barrière wordt dunner en $I$ groeit sneller dan lineair.

\textbf{9.} $\kappa' = \qty{0.89e10}{m^{-1}}$; bij \qty{0.5}{nm} zakt de
exponent van
$10.9$ tot $8.9$, $I \times 7$: de punt trekt zich
$\ln 7/2\kappa' \approx \qty{0.1}{nm}$ terug
--- een bult die elektronisch is, niet meetkundig.

\textbf{10.} $V(R_{\text{n}}) = 2 \times 90 \times 1.44/8 = \qty{32}{MeV}$;
$b = \qty{62}{fm}$.

\textbf{11.} $v = \qty{1.4e7}{m/s}$; $9 \times 10^{20}$ botsingen per seconde.

\textbf{12.} $G = 96$; $T = \eu^{-96} \approx 10^{-42}$.

\textbf{13.} Tempo $\sim 10^{-21}\,\unit{s^{-1}}$, halveringstijd
$\sim 10^{13}$ jaar --- duizend
keer de gemeten waarde: voor een exponent van honderd doet een
grof model het goed als het binnen enkele ordes van grootte landt.

\textbf{14.} $G = 79$: $\eu^{17} \approx 10^7$ keer korter, ongeveer $10^6$ jaar ---
één extra MeV, zeven ordes van grootte.

\textbf{15.} $G = 92$: $\eu^{4.5} \approx 90$ keer korter --- één fermi straal
is twee ordes van grootte.

\textbf{16.} Een proton is niet met positieve energie voorgevormd (het is
met zo'n \qty{7}{MeV} gebonden); een $^{12}$C heeft drie keer de lading en een
exponent die drie keer groter is.

\textbf{17.} Ja: het $\alpha$-deeltje heeft altijd \qty{4.2}{MeV}; zijn \omterm{def:b2:schrodinger-wave-functions:psi}{golffunctie}
heeft alleen een evanescent deel onder de barrière, waar geen meting
het met negatieve kinetische energie aantreft.

\textbf{18.} \qty{23.7}{GHz}; \qty{1.27}{cm}.

\textbf{19.} $(|{\uparrow}\rangle \pm |{\downarrow}\rangle)/\sqrt2$; de
symmetrische ligt lager --- geen knoop, minder kromming.

\textbf{20.}
$\psi = (\varphi_+\eu^{-\iu E_+t/\hbar} + \varphi_-\eu^{-\iu E_-t/\hbar})/\sqrt2$;
$P_{\downarrow} = \sin^2(\pi\nu t)$; periode \qty{42}{ps}.

\textbf{21.} $\eu^{-\Delta/k_BT} = 0.996$: vrijwel gelijk, geen inversie; een
inhomogeen elektrisch veld sorteert de bundel en stuurt de moleculen in de
bovenste toestand de holte in.

\textbf{22.} \qty{1.6e-23}{J}; $6 \times 10^{12}$ moleculen per seconde.

\textbf{23.} $\kappa a = 6.5$.

\textbf{24.} Zij wordt door het molecuul alleen vastgelegd, is voor elk
molecuul dezelfde en ongevoelig voor de buitenwereld; verschuivingen komen van het
\omterm{thm:b2:sound-waves:doppler}{dopplereffect} in de bundel, zwerfvelden, botsingen en
het trekken van de holte.

\textbf{25.} Opsluiting kwantiseert ($E_n \propto n^2/L^2$); evanescentie laat
de golf door ($T \sim \eu^{-2\kappa a}$); de exponenten worden bepaald door
$\sqrt{2m(V_0 - E)}/\hbar$ en de breedte --- massa, hoogte, afstand.
\end{solution}
